Ngưỡng cuối đời $\tau$ không lấy chung 0,8 như quy ước xe điện ở \Cref{sec:bai_toan}, vì nhiều cell trong các thí nghiệm này dừng thử trước khi SOH xuống 0,8. Tỉ lệ cell có SOH đo chạm 0,8 chỉ là 36\,\% ở XJTU, 77\,\% ở TJU, 0\,\% ở MIT và 8\,\% ở HUST, nên RUL theo ngưỡng 0,8 không xác định được cho phần lớn cell. Với mỗi bộ, $\tau$ được chọn trước khi học mô hình theo một quy tắc chỉ dựa trên phạm vi dữ liệu: mức thấp nhất trên lưới bước 0,01 mà ít nhất 80\,\% số cell của bộ có SOH đo chạm tới. Quy tắc này cho $\tau$ = 0{,}84 (XJTU), 0{,}81 (TJU), 0{,}87 (MIT) và 0{,}81 (HUST). Quy tắc dùng SOH đo của mọi cell trong bộ, kể cả cell kiểm tra, nhưng chỉ để định nghĩa mục tiêu RUL chung cho mọi mô hình, không thay đổi đầu vào hay tham số của mô hình nào. Cột ``cell có chu kỳ cuối đời'' của \Cref{tab:fcset} đếm theo định nghĩa chặt hơn dùng khi chấm, tức trung vị trượt 9 chu kỳ của SOH đo xuống dưới $\tau$, nên có bộ thấp hơn 80\,\% một chút. Vì $\tau$ khác nhau giữa các bộ, RUL của các bộ không so trực tiếp với nhau.
